\section{Signatures and absolute and relative discriminants}
\label{sec:global}

The subgroup lattice gives five transitive coset characters.  Evaluating
those characters on complex conjugation and on the released ramification
filtrations determines the signatures and conductor-discriminant exponents.
This section gives the complete derivation; the prime-by-prime ledger is
repeated in \cref{app:arithmetic-ledgers}.

\subsection{Archimedean signatures}

Let $E=K^H$ have degree $n=[G:H]$, and let $o_c(H)$ be the number of orbits
of a complex-conjugation subgroup on $G/H$.  Real embeddings are fixed
points and nonreal embeddings occur in transposed pairs, so
\[
 r_1+2r_2=n,\qquad r_1+r_2=o_c(H).
\]
Therefore
\begin{equation}
 r_1=2o_c(H)-n,\qquad r_2=n-o_c(H).
 \label{eq:signature-formula}
\end{equation}
The exact pairs $(n,o_c)$ are
\[
 (160,88),\quad(320,168),\quad(320,160),\quad(640,320)
\]
for $M$, $\Fplus/\Fthree$, $\Fzero$, and $L$.  Substitution in
\eqref{eq:signature-formula} gives
\begin{equation}
 (16,72),\quad(16,152),\quad(0,160),\quad(0,320).
 \label{eq:signatures}
\end{equation}
Each $r_2$ is even, so the sign $(-1)^{r_2}$ of every field discriminant is
positive.

\subsection{Conductor exponents}

The released lower filtration uses the ordered subgroup types
\[
 (I_3,P_3,Q_3,I_5,P_5,C_3,C_2,c_\infty).
\]
For a degree-$n$ coset carrier, write $o_U$ for the number of $U$-orbits.
For a local lower filtration $G_0\supseteq G_1\supseteq\cdots$, the Artin
conductor of the permutation character is
\begin{equation}
 a=\sum_{u\ge0}\frac{|G_u|}{|G_0|}\bigl(n-o_{G_u}\bigr).
 \label{eq:general-conductor}
\end{equation}
At $3$, the inertia group has order $18$, the wild group has order $9$, and
the order-$3$ group $Q_3$ occurs in six positive lower grades.  Their total
weights in \eqref{eq:general-conductor} are therefore
$1$, $9/18=1/2$, and $6(3/18)=1$.  At $5$, inertia has order $20$ and the
order-$5$ wild group occurs in three positive grades, giving weights
$1$ and $3(5/20)=3/4$.  A tame cyclic inertia group contributes only grade
zero.  Consequently the conductor-discriminant formula gives
\begin{align}
 v_3&=(n-o_{I_3})+\frac12(n-o_{P_3})+(n-o_{Q_3}),
 \label{eq:v3-formula}\\
 v_5&=(n-o_{I_5})+\frac34(n-o_{P_5}),
 \label{eq:v5-formula}\\
 v_A&=n-o_{C_3},\qquad v_B=n-o_{C_2}.
 \label{eq:vAB-formula}
\end{align}
Here $v_A$ is the common exponent at the three tame $C_3$ primes
$181,997,2346241$, and $v_B$ is the common exponent at the three reflection
$C_2$ primes $283,1801,q$.

The orbit counts and resulting exponents appear in
\cref{tab:global-orbits}.
\begin{table}[t]
\centering
\caption{Complete coset-orbit input for signatures and discriminants.  The
orbit vector is ordered as
$(I_3,P_3,Q_3,I_5,P_5,C_3,C_2,c_\infty)$.}
\label{tab:global-orbits}
\small
\begin{tabular}{@{}lcc@{}}
\toprule
field & orbit-count vector & $(v_3,v_5,v_A,v_B)$\\
\midrule
$M$ & $(22,28,56,8,32,64,80,88)$ & $(308,248,96,80)$\\
$\Fplus,\Fthree$ & $(36,56,112,16,64,128,160,168)$ & $(624,496,192,160)$\\
$\Fzero$ & $(28,56,112,16,64,128,160,160)$ & $(632,496,192,160)$\\
$L$ & $(56,112,224,32,128,256,320,320)$ & $(1264,992,384,320)$\\
\bottomrule
\end{tabular}
\end{table}
For example, the exponent of $3$ for $M$ is
\[
 (160-22)+\frac12(160-28)+(160-56)=308,
\]
and the other entries follow by the same substitutions.  The complete
character values rather than only the displayed totals are bound in the
exact certificate.

Define
\begin{align}
 q&=14932047182473291995860108491583652133938007263719,
 \label{eq:q-definition}\\
 A&=181\cdot997\cdot2346241=423395612137,
 \label{eq:A-definition}\\
 B&=283\cdot1801\cdot q\\
  &=7610610604104534884325967676315830570581925356194091077.
 \label{eq:B-definition}
\end{align}

\begin{proposition}[Signed absolute discriminants]
\label{prop:absolute-discriminants}
The signed field discriminants are
\begin{align*}
 \Disc(M)&=+3^{308}5^{248}A^{96}B^{80},\\
 \Disc(\Fplus)=\Disc(\Fthree)
  &=+3^{624}5^{496}A^{192}B^{160},\\
 \Disc(\Fzero)&=+3^{632}5^{496}A^{192}B^{160},\\
 \Disc(L)&=+3^{1264}5^{992}A^{384}B^{320}.
\end{align*}
Each field has exact finite ramified support \eqref{eq:released-support}.
\end{proposition}

\begin{proof}
Equations \eqref{eq:v3-formula}--\eqref{eq:vAB-formula} and
\cref{tab:global-orbits} give the prime exponents.  The even values of $r_2$
in \eqref{eq:signatures} give the positive signs.  The normal closure is
unramified outside \eqref{eq:released-support}, so each fixed field is also
unramified outside it.  Conversely, every exponent in each displayed
factorization is positive; after expanding $A$ and $B$, all eight primes
occur.  This proves exactness of the support as well as the factorizations.
\end{proof}

\subsection{Relative discriminants over the common core field}

For a finite extension $E/M$, the discriminant tower identity is
\begin{equation}
 \Disc(E/\Q)=\Disc(M/\Q)^{[E:M]}
        \Norm_{M/\Q}(\dd_{E/M}).
 \label{eq:disc-tower}
\end{equation}
Subtracting twice the exponent vector of $M$ from the three quadratic rows
gives
\begin{align*}
 \Fplus/M &: (624,496,192,160)-2(308,248,96,80)=(8,0,0,0),\\
 \Fzero/M &: (632,496,192,160)-2(308,248,96,80)=(16,0,0,0),\\
 \Fthree/M &: (624,496,192,160)-2(308,248,96,80)=(8,0,0,0).
\end{align*}
For the degree-four top,
\[
 (1264,992,384,320)-4(308,248,96,80)=(32,0,0,0).
\]
We have proved
\begin{equation}
\begin{aligned}
 \Norm_{M/\Q}\dd_{\Fplus/M}&=3^8,&
 \Norm_{M/\Q}\dd_{\Fzero/M}&=3^{16},\\
 \Norm_{M/\Q}\dd_{\Fthree/M}&=3^8,&
 \Norm_{M/\Q}\dd_{L/M}&=3^{32}.
\end{aligned}
\label{eq:relative-norms}
\end{equation}
In particular, all four relative extensions are unramified away from primes
of $M$ above $3$.  The absolute identity alone does not specify the local
factorization above $3$; that information comes from the complete
double-coset tables in the next section.
